\documentclass[11pt]{article}
\usepackage{fontspec}
\usepackage{amsmath,amssymb}
\usepackage[margin=1in]{geometry}
\usepackage{url}
\begin{document}
\sloppy
\section*{GROK H1411 (2026-10-06 14:11 PT): seed-probe budget calibration}
\textbf{P versus NP: OPEN.}

\section*{0. What landed since H1359}
\texttt{main} tip is \texttt{3b0812d1} (H1359). No peer commits under \path{research/agents/grok/} or \path{research/main/} since then.

\section*{1. Literature}
None. The next claim concerns this loop's own seed probe.

\section*{2. Budget calibration on the control seed (COMPUTATIONAL)}
H1359's probe (\path{seedprobe.py}, sha256 \texttt{79d2559b}, copied unchanged to \path{/workspace/scratch_h1411_1006}) decided nothing at 50\,000 conflicts, including the control seed $(0,4,14)$. This pass reran the control alone at three conflict budgets, in parallel, starting 14:11:52\,PT:
\begin{itemize}
\item 200\,000 conflicts: UNKNOWN after 15.9\,s.
\item 1\,000\,000 conflicts: SAT after 55.2\,s.
\item 4\,000\,000 conflicts: SAT after 57.4\,s, same model.
\end{itemize}
The returned eight-gate circuit is $(6,0,4,14)$, $(7,3,6,6)$, $(8,2,6,8)$, $(9,0,5,6)$, $(10,1,8,14)$, $(11,9,10,8)$, $(12,7,11,14)$, $(13,8,12,6)$; the script asserts that it computes $f=\texttt{a40ff54e580f5d0a}$ and starts with the seed, and its two-pebble BFS gives length $L=11$, re-simulated OK. This agrees with H1344 (every circuit found with this first gate pebbles in 11). So the control becomes decidable somewhere between $2\cdot10^5$ and $10^6$ conflicts. One satisfiable control says nothing about how many conflicts an UNSAT seed needs, so UNKNOWN at $10^6$ still does not mean UNSAT.

\section*{3. Rerun of the 74 seeds at $10^6$ conflicts (COMPUTATIONAL, partial)}
Launched 14:13:42\,PT, 2 workers (inputs \path{in0.txt} sha256 \texttt{ffb53524}, \path{in1.txt} \texttt{49d6b402}, 37 seeds each, control excluded). At 14:15\,PT: $(0,1,2)$ UNKNOWN in 89.6\,s and $(0,1,4)$ UNKNOWN in 88.4\,s; 72 still running into \path{run1m_0.out} and \path{run1m_1.out}. Nothing is concluded about other first gates yet.

\section*{4. Pending run read: tight $L=10$ (COMPUTATIONAL, partial)}
\path{l10.out} (H1258, 900\,s cap per cube) at 14:15\,PT: cubes 0 to 29 hit the cap; 0 SAT, 0 UNSAT, 30 timeouts out of 75 kept cubes. Timeouts are not UNSAT. $10\le L(f,2)\le 11$ is unchanged; whether $L(f,2)$ is 10 or 11 is OPEN.

\section*{5. Standing status}
No claim is falsified and no charged resource improves this pass. $10\le L(f,2)\le 11$ stands, the lower bound CONDITIONAL on the tightness lemma (CANDIDATE), the upper bound COMPUTATIONAL. The largest unconditional confirmed price stays one. The time-state-representation surface stays CANDIDATE. The SAT-to-charged-model compiler and the product-span obstruction are both OPEN. The D(4,8) semantic-cut bound of 48 stays FALSIFIED/RETIRED.

\section*{Next smallest step}
Read \path{run1m_*.out}: re-simulate any SAT seed independently and check whether its pebbling length is below 11. If most seeds stay UNKNOWN at $10^6$, stop raising budgets per seed and instead split the unseeded enumeration by the input pair of gate 6 (15 cubes) and run those as one long background job. Keep reading \path{l10.out}.

\bigskip
\noindent\textbf{Status sentence.} P versus NP remains OPEN: the SAT-to-charged-model compiler and the product-span obstruction are open load-bearing arrows. This pass finds that the seed probe decides the control seed $(0,4,14)$ SAT at $10^6$ conflicts (UNKNOWN at $2\cdot10^5$), returning an eight-gate circuit that pebbles with two registers in 11 steps, that the first two of 74 other seeds stay UNKNOWN at $10^6$, and that tight $L=10$ cubes 0 to 29 time out (0 SAT, 0 UNSAT); $10\le L(f,2)\le 11$ is unchanged, with the lower bound CONDITIONAL on the tightness lemma. No claim is falsified, and this is not a proof of anything about P versus NP.
\end{document}
